\section{The complete parent and its positive remainders}\label{app:parent}
For a direct expansion of \eqref{eq:slack}, put
\[
 q=v\zeta-Tuw=-T\sqrt x\,\Qh,\qquad
 m_p=\frac{\mu p}{1-p},
\]
and use the two independent normals $N,M$ in \eqref{eq:normal}.
Let
\begin{align}
 F_s={}&q_c\Qh^2+nj^2y\Bh^2+r_c\Rh^2+d_cV^2\Qh^2
 +(np^2y^2+\half p^2y^2-g_c)V^2\Bh^2\notag\\
 &+f_c\Ps+g_cV^2\Ps+k_c\Zc^2(V^2-1)^2,\notag\\
 L_N={}&nj\bc+npv^2\bc-j\alpha\bc+\half pv^2\bc
                         +K_RR_u+(K_c/j)uvq,\notag\\
 L_M={}&C_qq+\mu uv\bc+v(\alpha-BTCHu^2)R,\notag\\
 H_N={}&h_0(1+u^2)+\half v^2+\alpha u^2+n(1+v^2),\notag\\
 H_M={}&(1+u^2)[C(1+v^2)+D+\half]+A(1+v^2)+(\alpha+\half)v^2 .
 \label{eq:parentcoeff}
\end{align}
The full parent is the polynomial identity
\begin{equation}
 D_vS_c=F_s+2L_NN+2L_MM+H_NN^2+H_MM^2+
              2(\alpha+\half)uvNM+m_p(1+u^2)N^2 .
 \label{eq:fullparent}
\end{equation}
It follows by inserting \eqref{eq:rotation} into \eqref{eq:slack} and
using \eqref{eq:polys}--\eqref{eq:source}.
Equivalently these are the homogeneous RFN8/RFN12 identities of the
preserved companion source \texttt{latex/certificates/c2\_short.tex}.
The four scalar cancellations needed in this expansion are, with
$K_d=CH-(D+\half)(1-H)$,
\begin{align*}
 0&=T(K_d+AH+\half)-A,\\
 0&=jK_d+\mu+xTCH,\\
 0&=jK_d+xTCH+T[CH+AH+(\alpha+\half)H]
                   -A-(\alpha+\half)(1-p)+\alpha/B,\\
 0&=np-\mu+p/2+[-\mu-(\alpha+\half)(1-p)]x+j(\alpha-n)+yK_c .
\end{align*}
They are exact rational identities in the printed source dictionary.

The retained positive term in \eqref{eq:coreidentity} is explicitly
\begin{align}
 \calP={}&m_p(1+u^2)N^2+(h_0+\alpha)(uN)^2
             +\half(vN)^2+n[v(N+p\bc)]^2\notag\\
 &+(C+D+\half)(uM)^2+(C+A+\alpha+\half)(vM)^2+C(uvM)^2\notag\\
 &+d_c(V\Qh)^2+\half p^2v^2\bc^2
            +g_c[V(w-UV\Zc)]^2+\kappa(V\Qh-\Rh)^2
            +g_c(\Qh^2+\Rh^2).
 \label{eq:retainedall}
\end{align}
In particular $d_c=H^2(A+C+\alpha+\half)y^2q_z+\kappa>0$.
Combining
\[
 h_0N^2+n(N+j\bc)^2=\widehat N^2+E_*^2,\quad
 g_cV^2(\Ps-\Bh^2)=g_c(\E-\Ps),\quad
 -\kappa(V\Qh+\Rh)^2
 =\kappa(V\Qh-\Rh)^2-2\kappa V^2\Qh^2-2\kappa\Rh^2
\]
with \eqref{eq:fullparent} gives precisely \eqref{eq:coreidentity}
and the complete cross terms \eqref{eq:cross}.
Thus the decomposition is an identity for all four state coordinates.

\subsection*{The two radial sums of squares}
The radial polynomials used in \eqref{eq:radialpurchase} are
\begin{align}
 \W&=4\Ps+4U^2\Qh^2+4U^2\Rh^2-V^2(U^2-1)^2w^2,\notag\\
 \W^\pi&=4\Ps+4U^2\Qh^2+4U^2\Bh^2-\Zc^2(U^2-1)^2 .
 \label{eq:radialpolys}
\end{align}
Here is their complete positive Gram representation on the chart
$U,V\geq0$, $w\Zc\geq0$.
Put
\[
 t_7=\begin{pmatrix}
 \Zc-U^2Vw\\Vw-U^2Vw\\Uw-U^2Vw\\U\Zc-U^2Vw\\
 UVw-U^2Vw\\UV\Zc-U^2Vw\\U^2w-U^2Vw
 \end{pmatrix},\qquad
 t_3=\begin{pmatrix}1-UV\\V-UV\\U-UV\end{pmatrix}
\]
and
\[
 t_7^\pi=\begin{pmatrix}
 Vw-U^2\Zc\\\Zc-U^2\Zc\\Uw-U^2\Zc\\UVw-U^2\Zc\\
 U\Zc-U^2\Zc\\UV\Zc-U^2\Zc\\U^2w-U^2\Zc
 \end{pmatrix},\qquad
 t_3^\pi=\begin{pmatrix}w-U\Zc\\\Zc-U\Zc\\Uw-U\Zc\end{pmatrix}.
\]
With the following two fixed rational matrices,
\[
 G_7=\begin{pmatrix}
4&-2&-3/2&0&-1/2&0&1/2\\
-2&3&0&-3/2&0&-1&2\\
-3/2&0&4&3/2&-2&-3&0\\
0&-3/2&3/2&4&-3/2&0&-3/2\\
-1/2&0&-2&-3/2&3&2&-1\\
0&-1&-3&0&2&8&-4\\
1/2&2&0&-3/2&-1&-4&4
\end{pmatrix},\qquad
 G_3=\begin{pmatrix}3&-2&-2\\-2&2&2\\-2&2&3\end{pmatrix},
\]
the identities are
\begin{align}
 \W&=t_7^TG_7t_7+4Vw\Zc(1-U)^2+Uw\Zc\,t_3^TG_3t_3,\notag\\
 \W^\pi&=(t_7^\pi)^TG_7t_7^\pi+4Vw\Zc(1-U)^2
                                      +UV(t_3^\pi)^TG_3t_3^\pi .
 \label{eq:W19}
\end{align}
Their leading principal minors are
$(4,8,101/4,1049/16,907/16,3279/16,743/16)$ and $(3,2,2)$.
Both matrices are therefore positive definite. These explicit polynomial
identities are the homogeneous RFN19 and permuted RFN19 identities.
All extra square roots have fixed testing arguments.

For $w\Zc\leq0$, one instead uses the following direct expansions:
\begin{align}
 \W={}&(4U^2+3V^2+4U^4+3U^4V^2+2U^2V^2)w^2
 +(4+4U^2+8U^2V^2)\Zc^2
             -16UV(1+U^2)w\Zc,\notag\\
 \W^\pi={}&(4U^2+4V^2+4U^4+4U^2V^2)w^2
 +(3+2U^2+3U^4+8U^2V^2)\Zc^2
             -16UV(1+U^2)w\Zc .
 \label{eq:negativeW}
\end{align}
Every displayed monomial is a nonnegative coefficient times a square,
using the fixed radical $\sqrt{-UVw\Zc}$ for the last terms.
Together with the five-factor expansion of $\Ps$ in the main text, these
give a direct sum-of-squares identity also on the negative state chart.

\subsection*{The state signs and testing axes}
Let $K_d^{\mathrm{test}}$ denote the $d=1$ core in the numerator
of \eqref{eq:slack}, put $h=1+u^2+v^2$,
$\Delta=h^2-u^2v^2>0$, and set
\[
 A_0=\half\begin{pmatrix}h&-uv\\-uv&h\end{pmatrix},\qquad
 B_0=\half\begin{pmatrix}-u^2v&-uv^2\\uv^2&u^2v\end{pmatrix}.
\]
The exact core completion is
\begin{align}
 K_d^{\mathrm{test}}/D_u
 ={}&\left[\binom tr+A_0^{-1}B_0\binom w\zeta\right]^TA_0
 \left[\binom tr+A_0^{-1}B_0\binom w\zeta\right]\notag\\
 &+\frac{D_uD_v}{2\Delta}
      [u^2v^2(w^2+\zeta^2)-2uvh\,w\zeta] .
 \label{eq:signcore}
\end{align}
Since $A_0\succ0$, a negative core requires $uvw\zeta>0$.
Each of the nine original source summands is invariant under the gauges
\[
 (u,v,W,X,Y,Z)\longmapsto
 (\varepsilon u,\eta v,W,\eta X,\varepsilon Y,\varepsilon\eta Z).
\]
Thus a negative core can always be reduced to $u,v>0,w\zeta>0$.
At $u=0,v=0,w=0$ or $\zeta=0$, the remainder in \eqref{eq:signcore}
is nonnegative. The explicit identities on both state charts extend to these axes
with their literal zero radicals, establishing the theorem on the closed cap.
